\cotefile{111}
\scanpage{2}\mspage{\handwritten{\uline{1}}}\margin{\handwritten{notes du \uline{10.9.68}}}

Tapis de Quillen.

I \uline{Relations entre catégories et ensembles semi-simplic\unsure{aux|iaux}}.

A toute catégorie $C$, on associe un ensemble semi-simpli\deleted{\unsure{v}}cial $S(C)$, trouvant \unsure{ninsi|ainsi} un foncteur pleinement fidèle \deleted{de Cat}
\[
S;\ (\mathrm{Cat})\handwritten{\longrightarrow}(\mathrm{Ssimpl})\ .
\]

\deleted{Si $C$ est associé à} Les systèmes locaux d'ensemble sur $SC$ correspondent aux foncteurs sur $C$ qui transforment toute flêche en isomorphisme (i.e. qu\deleted{\unsure{e}}i se factorisent par le groupoïde associé à $C$). Les $H^i$ sur $SC$ d'un tel système local ($H^{\unsure{o|0}}$ pour ens, $H^1$ pour groupes, $H^i$ quelconques pour groupes abéliens) s'interprètent en termes des foncteurs $\underset{\handwritten{\leftarrow}}{\mathrm{lim}}^{(i)}$ déri\deleted{\unsure{v}}gés de $\underset{\handwritten{\leftarrow}}{\mathrm{lim}}$, ou si on préfère, des $H^i($ du \uline{topos} $C^{\handwritten{\wedge}}$. On voit ainsi à quelle condition un foncteur $C\handwritten{\to}C'$ induit un\deleted{\unsure{e}} homotopisme $SC\handwritten{\to}SC'$ : En vertu du critère cohomologique de Artin-Mazur, il f et s que pour tout système de coéfficients $F'$ sur $C'$, l'homomorphisme naturel $\underset{C'}{\underset{\handwritten{\leftarrow}}{\mathrm{lim}}}^{(i)}F'\handwritten{\longrightarrow}\underset{C}{\underset{\handwritten{\leftarrow}}{\mathrm{lim}}}^{(i)}F$ soit un isomorphisme (pour les $i$ pour lesquels cela a un sens). \deleted{En termes}

A $C$ on peut asso\deleted{\unsure{v}}c\unsure{er|ier|ér} le topos $\handwritten{\widetilde{C}}$, qui varie de façon \uline{covariante} avec $C$. (NB le foncteur $C\handwritten{\mapsto\widetilde{C}}$ n'a plus rien de pleinement fidèle, semble-til ??) Les systèmes de coéfficients \inserted{ensemblistes} sur $C$ ($\overset{\mathrm{déf}}{=}$ les foncteurs \inserted{$C^{o}\handwritten{\to}\mathrm{Ens}$} transformant isomorphismes en isomorphismes) correspondent aux faisceaux localement constants i.e. les objets localement constants de $\handwritten{\widetilde{C}}$, définis intrinsèquement en \unsure{termes} de $\handwritten{\widetilde{C}}$. Ainsi, le fait pour un foncteur $F:C\handwritten{\to}C'$ d'induire un\deleted{e équivalence} homotopisme $S(C)\handwritten{\to}S(C')$ ne dépend que du morphisme de topos $\handwritten{\widetilde{F}}:\handwritten{\widetilde{C}}\handwritten{\to}\handwritten{\widetilde{C}}'$ induit, et signifie que pour tout faisceau localement constant $F'$ sur $C'$ \unsure{i.}e. sur $\handwritten{\widetilde{C}}'$, les applications induites $H^i(\handwritten{\widetilde{C}}',F')\handwritten{\longrightarrow}H^i(\handwritten{\widetilde{C}},\handwritten{\widetilde{F}}^{\handwritten{*}}(F'))$ sont des isomorphismes (pour les $i$ pour lesquels cela a un\deleted{\unsure{es|e}} sens).

% archive sheet number: 1
\scanpage{3}\mspage{\handwritten{2}}

On a aussi un foncteur évident
\[
T:(\mathrm{Ssimpl})\handwritten{\longrightarrow}(\mathrm{Cat})\ ,
\]
en associant à tout ensemble semi-simplicial $X$ la catégorie $T(X)=\handwritten{\Delta}_{/X}$ des simplexes sur $X$, dont l'ensemble des objets est la réunion disjointe des $X_n$ .... (c'est une catégorie fibrée sur la catégorie $\handwritten{\Delta}$ des simplexes types, à fibres les catégories discrètes définies par les $X_n$). Ceci posé, Quillen prouve que pour tout $X$, $ST(X)$ est isomorphe canoniquement à $X$ \unsure{dans} la catégorie homotopique construite avec $(\mathrm{Ssimpl})$, et que pour toute $C$, la catégorie $TS(C)$ est \inserted{canoniquement} ``homotopiquement équivalente à $C$'' \deleted{dans le sens} i.e. \inserted{canoniquement} isomorphe à $C$ dans la catégorie quotient de $(\mathrm{Cat})$ obtenue en inversant les foncteurs qui sont des homotopismes. \deleted{Il en résulte} Ces isomorphismes sont fonctoriels en $X$. Il en résulte formellement qu'un morphisme $f:X\handwritten{\unsure{-|\to}}Y$ dans $(\mathrm{Ssimpl})$ est un homotopisme si et seulement si il en est ainsi de $T(f):\unsure{T(X)}\handwritten{\to}T(Y)$, d'où des foncteurs $S':(\mathrm{Cat})'\handwritten{\longrightarrow}(\mathrm{Ssimpl})'$ et $T':(\mathrm{Ssimpl})'\handwritten{\to}\unsure{(\mathrm{Cat})'}$ entre les catégories ``homotopiques'', \inserted{construites avec $(\mathrm{Cat})$ resp. $(\mathrm{Ssimpl})$,} qui sont quasi-inverses l'un de l'autre.

De plus, Quillen construit un isomorphisme canonique et fonctoriel dans $(\mathrm{Cat})'$ entre $C$ et la catégorie opposée $C^o$, ou ce qui revient au même, un isomorphisme canonique et fonctoriel dans $(\mathrm{Ssimpl})'$ entre $S(C)$ et $S(C^o)$. \deleted{Cela lui permet alors, pour} La définition est telle que le foncteur induit sur les systèmes locaux sur $C$ tra\unsure{hsfôrme|hsforme|nsforme} le foncteur contravariant $F$ sur $C$, transformant toute flêche en flêche inversible, en le foncteur covariant (i.e. contravariant sur $C^o$) ayant mêmes valeurs sur les objets, et obtenu sur les flêches en remplaçant $F(u)$ par $F(u)^{-1}$; en d'autres termes, l'effet de l'homotopisme de Quillen sur les groupoïdes fondamentaux est l'isomorphisme évident entre les groupoïdes fondamentaux de $C$ et de $C^o$, \deleted{\unsure{qui est l'id en}} compte tenu que le deuxième est l'opposé du

% archive sheet number: 2
\scanpage{4}\mspage{\handwritten{3}}

premier. Comme application, Quillen obtient une interprétation \unsure{faisceautique} de la cohomologie d'un ensemble semi-simplicial à coéfficients dans \unsure{un} système \inserted{\uline{local}} contravariant \deleted{\unsure{de}} \inserted{$F$} (défini classiquement par le complexe \unsure{cosimplicial} des $C^n(F)=\handwritten{\coprod}_{x\in X_n}F(x)$)\margin{\handwritten{$\prod$ ?}} : on considère le système local contravariant défini par $F$, on l'interprète comme un faisceau sur \inserted{$T(X)$ i.e. objet de} $(\mathrm{Ssimpl})_{/X}$, et on prend sa cohomologie. - Cependant, quand $F$ est un système de coéfficien\unsure{bs|ts|s} covariant pas nécessairement local, on n'a tjrs pas d'interp\unsure{éé|ré}tation de \unsure{ses} groupes de cohomologie classiques en termes faisceautiques; ni, lorsque \unsure{$F$} est contravariant, de son homologie, ou inversement de sa cohomologie faisceautique en termes classiques.\margin{\handwritten{Si ! \tikz[scale=.18,baseline]{\draw (.5,1)--(0,1)--(0,-1)--(.5,-1);}}}

\margin{\handwritten{\(\Vert\)}}% parallel pencil rules beside the Quillen assertion
A propos de la notion de foncteur qui est un homotopisme. Quillen montre qu\deleted{e}'un tel foncteur $F:C\handwritten{\to}C'$ induit une équivalence entre la sous catégorie triangulée \inserted{$D^b_{lc}(C')$} de la catégorie \inserted{dérivée \inserted{bornée} de celle \deleted{des}} des faisceaux abéliens sur $C'$, dont les faisceaux de cohomologie sont des systèmes locaux, et la catégorie \unsure{analogue} pour $C$; \deleted{\unsure{il s'en}} et réciproquement. \deleted{Pour la partie directe} On peut dans cet énoncé introduire aussi n'importe quel anneau de base (à \unsure{condition} de le supposer $\handwritten{\ne}0$ dans le cas de la réciproque); \deleted{constant} la partie \unsure{directe} vaut aussi avec un anneau de boefficients pas néc. constant, mais \unsure{constant} tordu. Je pense que ce résultat (facile) doit pouvoir se généraliser \unsure{ainsi} Soit $f:X\quad X'$ un morphisme de topos qui soit tel que pour tout \deleted{\unsure{système}} faisceau loc. const. sur $X'$, $f$ induise un isomorphisme sur les \unsure{cohomologies} (avec cas non commutatif inclus). Supposons que \deleted{\unsure{le topos $X$|$X$ (resp. $X'$)}} \inserted{$X$\deleted{\unsure{$'$}} et $\handwritten{X'}$} soit \uline{\unsure{localement} homotopiquement trivial}, i.e. que \deleted{tout} pour tout entier $n\inserted{\handwritten{\geq}1}$, tout objet $U$ \deleted{\unsure{de $X'$}} ait un \deleted{revêtement}\inserted{recouvrement} par des $U_i\handwritten{\to}U$, tels que a) tout système local sur \unsure{$U$} devient constant sur $U_i$, et toute section sur $U$ devient constante sur $U_i$ et b) pour tout \deleted{système} groupe abélien $G$, les $H^{\deleted{\unsure{i}}j}(U,G)\handwritten{\to}H^j(U_i,G)$ sont nuls pour $1\handwritten{\leq}j\handwritten{\unsure{<|\leq}}n$.\margin{\handwritten{\tikz[scale=.16,baseline]{\draw (0,0) circle (.28) (.15,.2)--(1.25,1.75);}\ \(\vert\!\vert\!\vert\)}}% callout and three vertical strokes beside the condition
\margin{\handwritten{\tikz[scale=.16,baseline]{\draw (0,0) circle (.28) (.15,.2)--(1.25,1.75);}\ Attention, cette condition n'est typiquement pas satisfaite par les schémas en leur \unsure{topologie} \unsure{étale} (mais bien par \unsure{schémas} \unsure{noeth.|loc. noeth.} avec top. Zariski).}} Alors le foncteur $D^b_{lc}(X')\handwritten{\to}D^b_{lc}(X)$ induit par $f$ est

% archive sheet number: 3
\scanpage{5}\mspage{\handwritten{4}}

\deleted{pleinement fidèle (resp.} une équivalence). Même énoncé si on met dans le coup un \deleted{Anneau \unsure{loc.}} système local d'anneaux sur $X'$. Ehfin, $f$ induit \deleted{une} \deleted{isomorphisme}\inserted{\deleted{foncteur pleinement fidèle (resp.} une équivalence)} entre la catégorie des coefficients locaux sur $C$ et celle des coéff. locaux sur $C'$ \deleted{(?? peut-être \inserted{cette propriété doit} \unsure{cela} doit-il être mis dans les hypothèses préliminaires sur $f$ ?).} Principe de démonstration: on commence par prouver ce dernier résultat, en notant que \deleted{dans le ca} si un topos est localement hom. trivial, il est loca connexe et loc. simplement connexe, d'où une bonne théorie du $\handwritten{\pi}_{\unsure{o|0}}$ et du $\handwritten{\pi}_1$ (qui sont ici discrets), et on est ramené \deleted{pour} à un cas particulier du critère d'homotopisme de Artin et Mazur, savoir un critère \inserted{cohomologique} pour qu'un homomorhisme de groupes $G\handwritten{\to}H$ (ici des groupes $\handwritten{\pi}_1$ de $X,X'$) soit un isomorphisme: il doit induire des isomorphismes sur les $H^{\unsure{o|0}}$ et $H^1$ (y inclus dans le cas non commutatif...). On prouve la pleine fidélité en se ramenant par la suite spectrale aux cas des $\mathrm{Ext}^i$ de coéfficients locaux. Par la suite spectrale encore, cela résultera du fait suivant: si $X^{\deleted{'}}$ est \unsure{homalement|honalement|localement} hom. trivial, alors la catégorie des faisceaux abéliens loc. constants est stable par $\underline{\mathrm{Ext}}^i$, et le foncteur $M\handwritten{\mapsto}M_X$ des \deleted{faisceaux constants} $(\mathrm{Ab})$ dans $X_{ab}$ commute aux dits $\mathrm{Ext}^i$.

En fait, $X$ et $X'$ étant loc.hom\deleted{p}triviaux, les conditions suivantes sur $f$ seront équivalentes:

a) $f$ est un homotopisme, i.e. induit pour tout système local (pas néc commutatif) sur $X'$ un isomorphisme sur les $H^i$\deleted{pour \unsure{$i<n$|$i\le n$}}.

b) \deleted{Même chose ave} $f$ induit une équivalence $D^b_{lc}(X')\handwritten{\longrightarrow}D^b_{lc}(X)$.

a\handwritten{'}) $f$ induit une équivalence entre la catégorie des systèmes locaux abéliens sur $X'$ et $X$, et des isomorphismes sur les $H^i$ correspondants (\unsure{donc} on ne prend ici que des coéfficients commutatifs).

\margin{\handwritten{?}} J'ignore si on peut dans a) se borner aux systèmes locaux commutatifs \deleted{i.e. dans}.

L'équivalence entre a) et b) fournit une première justification ou motivation pour définir des types d'homotopie via la catégorie $D\handwritten{^{b}_{lc}(X)}$,

% archive sheet number: 4
\scanpage{6}\mspage{\handwritten{5}}

éventuellement muni de la sous-catégorie pleine de tous les systèmes \unsure{locaux} sur $X$, et du foncteur cohomologique sur $D^b_{lc}(X)$ à valeurs dqns \unsure{laddtë|laddte} catégorie, et bien sûr du produit tensoriel (mais alors on sort de $D^b$ pour entrer dans $D^-$, redactor demerdetur).

II \quad $n$-catégories, catégories $n$-uples, et $(\mathrm{Gr})$-catégories.

Appelons bi-catégorie une catégorie dans$(\mathrm{Cat})$, i.e. un triple $(C^{\deleted{o}}_{\handwritten{0}},C^{\deleted{1}}_{\handwritten{1}},\ \deleted{\unsure{)\text{ où }}}p_1,p_2,\handwritten{\pi})$, où $C^o$ est une catégorie, $p_1,p_2:C^{\deleted{1}}_{\handwritten{1}}\handwritten{\rightrightarrows}C^{\deleted{o}}_{\handwritten{0}}$ deux foncteurs, $\handwritten{\pi}:(C^{\deleted{1}}_{\handwritten{1}},p_1)\mathbin{x}_{C^o}(C^{\deleted{1}}_{\handwritten{1}},p_2)\handwritten{\to}C^{\deleted{1}}_{\handwritten{1}}$ un foncteur, avec les axiomes habituels. Interprétant $(\mathrm{Cat})$ comme la sous-catégorie pleine de $(\mathrm{Ssimpl})$ des objets qui .... , et de même pour les catégories à valeurs dans une catégorie telle $(\mathrm{Cat})$, on trouve que la catégorie \inserted{$(\mathrm{Bicat})$} des $(\mathrm{Cat})$-catégories est équivalente à la sous-catégorie pleine de la catégorie des ensembles bi-semisimpliciaux

\diagram{\(\begin{array}{cccc}
X^{\deleted{oo}}_{\handwritten{00}} & X^{\deleted{o1}}_{\handwritten{01}} & \ldots & X^{\deleted{o,n}}_{\handwritten{0n}}\\
X^{\deleted{1o}}_{\handwritten{10}} & X^{\deleted{11}}_{\handwritten{11}} & \ldots & X^{\deleted{1,n}}_{\handwritten{1n}}\\
\multicolumn{2}{l}{.......} & &
\end{array}\)}

formée des objets qui ligne par ligne et colonne par colonne \unsure{satisfoht|satisfont} la condition d'exctitude familière (transformer sommes amalgamées en produits fibrés). De même la catégorie \deleted{des} $(\mathrm{Tricat})$ des $(\mathrm{Bicat})$-catégories \unsure{ssinterprète|s'interprète} en termes d'objets tri-\inserted{semi}simpliciaux, etc.

Exemple: soit $C$ une catégorie, posons $C^{\deleted{o}}_{\handwritten{0}}=C$, $C^{\deleted{1}}_{\handwritten{1}}=\underline{\mathrm{Fl}}(C)$, avec structures $p_1,p_2$, évidentes (d'où une $(\mathrm{Cat})$-catétorie dont les composantes semi-simpliciales supérieures sont les catégories $\underline{\mathrm{Fl}}^n(C)$ des $n$-flêches de $C$ ...). L'ensemble bi-semisimplicial correspondant a pour composantes les \unsure{ensembles} $\mathrm{Hom}(\handwritten{\Delta}^n\mathbin{x}\handwritten{\Delta}^m,C)$. Par exemple, les objets de $X^{\deleted{11}}_{\handwritten{11}}$ sont les carrés commutatifs de flêches de $C$, qu'on peut interpréter (de deux manières différentes) \unsure{comme}

% archive sheet number: 5
\scanpage{7}\mspage{\handwritten{6}}

morphismes de flêches.

Notons une symmétrie évidente dans la catégorie $(\mathrm{Bicat})$, interprétée en termes d\deleted{e}'ensemble bi-semisimpliciaux, obtenue en échangeant les \unsure{lignes} et les colonnes dqns \deleted{les} de tels ensembles bi... Un ensemble \unsure{bisemisimplicial} qui ... peut donc de deux façons différentes être considéré comme définissant une bicatégorie. De même le groupe symétrique $\mathrm{Sym}_n$ opère sur la catégorie \deleted{de} $(n\unsure{-|{}}\mathrm{upcat})$ des catégories $n$-uples ..

Une 2-catégorie peut être définie comme \unsure{constttuée|constituée} par un ensemble \deleted{discret} $C^{\deleted{o}}_{\handwritten{0}}$ (l'ensemble des objets ), qu'on interprètera comme une catégorie discrète, plus la donnée pour tout $x,y\handwritten{\in}C^{\deleted{o}}_{\handwritten{0}}$ d'une catégorie $\underline{\mathrm{Hom}}(x,y)$, plus \deleted{une} \inserted{des foncteurs} composition entre celles-ci ... Prenant la catégorie somme \deleted{de} $C^{\deleted{1}}_{\handwritten{1}}$ des $\underline{\mathrm{Hom}}(x,y)$, On peut exprimer la donnée de ces foncteurs de composition, et des sources et buts \inserted{$x,y$} dans les catégories $\underline{\mathrm{Hom}}(x,y)$, comme définis par des foncteurs ``source'' et ``but'' $C^{\deleted{1}}_{\handwritten{1}}\handwritten{\to}C^{\deleted{o}}_{\handwritten{0}}$, et un foncteur de composition $C^{\deleted{1}}_{\handwritten{1}}\mathbin{x}_{C^{\deleted{o}}_{\handwritten{0}}}C^{\deleted{1}}_{\handwritten{1}}\handwritten{\to}C^{\deleted{1}}_{\handwritten{1}}$, les axiomes des 2-catégories correspondant exactement aux axiomes des$(\mathrm{Cat})$-catégories. De cette façon, la catégorie des \unsure{2-catégories} est équivalente à la catégorie des $(\mathrm{Cat})$-catégories qui ont une ``\unsure{catégorie} d'objets'' qui est discrète, ou encore à la catégorie des Bicatégories qui, interprétées comme ensembles bi-semisimpliciaux, ont \deleted{une} comme première ligne un ensemble \unsure{bemisimplicial|semisimplicial} discret (i.e. un foncteur constant sur $\handwritten{\Delta}$). Il parait qu'on a une interprétation simple dans le même esprit de la notion de $n$-catégorie en termes d'ensembles \unsure{$n$-semi-simpliciaux}.

\margin{\handwritten{\(\Vert\)}}% double marginal stroke beside Demander
\deleted{au} Demander à Quillen relations entre la notion de $m$-Catégorie dans la catégorie $(n\text{-}\mathrm{Cat})$, et celle d\deleted{e}\handwritten{'} ensemble $(m+n)$-semisimplicial. \deleted{(en particulier entre cette dernière notion et celle de}

L'interprétation standard semisimpliciale permet d'interpréter

% archive sheet number: 6
\scanpage{8}\mspage{\handwritten{7}}
% Two lower corners of a drawn frame enclose the note above the first typed line.
\margin{\handwritten{Attention : la catégorie \uline{Hom} \unsure{gr.ab.|grab}, qui s'identifie à la catégorie définie par un $\unsure{\mathbb{R}|\mathrm{R}}\mathrm{Hom}_{\unsure{\mathbb Z|Z}}(C_{\unsure{*|\bullet}},C'_{\unsure{*|\bullet}})$ ; ce point devrait être élucidé !}}

de même les $(\mathrm{Cat})$-groupes(ou groupes à valeurs dans $(\mathrm{Cat})$) comme étant les groupes semisimpliciaux qui ``sont''des catégories i.e. qui \deleted{se définissent} satisfont la condition d'exactitude familière, permettant d'en déterminer tous les composants à l'aide de $C_{o\deleted{\unsure{o}}}$ et $C_1$. Appliquant le tapis de Dold, dans le cas commutatif, on trouve que les $(\mathrm{Cat})$-groupes abéliens forment une catégorie équivalente à celle des complexes de chaines de groupes abéliens qui sont de longueur 1:
\[
0\handwritten{\leftarrow}C_o\handwritten{\leftarrow}C_1\handwritten{\leftarrow}0\ .
\]

Quand on explicite la $(\mathrm{Cat})$-groupe abélien associé à un tel complexe de chaines, on trouve ``ma'' construction chérie de la ``catégorie de Picard'' associée au complexe: objets ceux de $C_o$, flêches de $x$ dans $y$ les $z$ dans \unsure{$C_1$} tels que $dz=y-x$. On trouve donc à isomorphisme ppès tous les \unsure{groupes} abéliens de $(\mathrm{Cat})$ ainsi. Un complément important est que toute catégorie de Picard \deleted{(resp. et commutative)} est équivalente (au sens techn\unsure{iqque|ique} qu'on ne précise pas ici) à \deleted{la catégorie} un vrai $(\mathrm{Cat})$-groupe \deleted{(resp. $(\mathrm{Cat})$-groupe commutatif)}.
% slanted
\margin{\handwritten{Quillen dixit, mais cela \unsure{contredit} les résultats \unsure{de ma|du} \unsure{lettre|livre} \unsure{à} \unsure{Verdier|Notations}, \unsure{la|voir} \unsure{même|sur|fin} \unsure{que|la|de la} \illegible{1} ? Mais attention, \unsure{même|vrai} si on part d'une cat. de Picard commut., elle ne se \unsure{réalise} que par un $(\mathrm{Cat})$-gr. qui \uuline{non} \uline{\unsure{comm.|commut.}}}}
\margin{\handwritten{Quand on simplifie mod ``\unsure{Grab|Gr.ab}'' équivalence cela correspond à travailler dans la catégorie \uline{dérivée} des complexes, \unsure{seuls} invariants \unsure{sont} bien évidemment $\pi_0,\pi_1$).}}

On notera qu'on a une \deleted{\unsure{interp}} classification beaucoup plus \unsure{restreinte} que pour les H-espaces (resp. les H-espaces homotopiquement comm.) mod équivalence d'homotopie, de façon précise entre les groupes resp. les groupes abéliens de la catégorie homotopique. La raison en est qu'on est beau\deleted{\unsure{x|c}}coup plus restrictif dans le contexte catégorique pour le sens de \unsure{l'associativité} (resp. et de la commutativité), puisqu'on ne demande pas seulement l'existence d'isomorphismes d'associ\unsure{aitivité|ativité}\deleted{\unsure{,|$\times$}} (qui suffirait pour avoir un ``groupe dans \inserted{la catégorie homotopique} $(\mathrm{Cat})'\handwritten{\unsure{\cong|\simeq}}(\mathrm{Ssimpl})'$, modulo l'existence d'unités et d'\unsure{inverses} mais qu'on suppose \uline{donné\deleted{e}s} ces isomorphismes, de façon à satisfaire à certaines compatibilités: c'est ce qui permet, me semble-t-il, de prouver \unsure{qu'à} équivalence près, on peut supposer qu'il s'agit de vrais groupes dans $(\mathrm{Cat})$, \deleted{resp. de vraie commutativité.}
% slanted
\margin{\handwritten{\deleted{\unsure{pourquoi|prouver}}\ \(\bigl(\)}}% opening bracket beside the final two lines

% archive sheet number: 7
\scanpage{9}\mspage{\handwritten{8}}

On doit trouver une approximation meilleure d'un cran à la classification topologique en prenant des $(2\text{-}\mathrm{Cat})$-groupes resp. des $(2\text{-}\mathrm{Cat})$-\unsure{groupes} abéliens, ou plus généralement des `` Gr-2-catégories '' resp. ``Grab-2-catégories'' . (NB U\deleted{\unsure{p}}ne Grab-catégorie défi\deleted{\unsure{b}}nit une Grab-\unsure{2catégorie} comme une catégorie définit une 2-catégorie; du point de vue semi-simplicial, cette opération serait analogue à celle de ``cosquélette\deleted{\unsure{x|r|\({}^{\times}\)}} tronqué''; Mais également une Grab-2-catégorie défini\deleted{e}\handwritten{t} \deleted{en}\handwritten{une} Grab-catégorie\deleted{\unsure{x|\({}^{\times}\)}} , savoir la catégorie des automorphismes de l'objet unité; cette \deleted{dé} opération serait analogue à celle de ``tuage des groupes d'homotopie en dim $\handwritten{>1}$''). Pour la notion gro\deleted{\unsure{t}}ssière ( Gr\deleted{ab} ou Grab) l'idée de ma lettre à Verdier était \unsure{qu'on} doit trouver la même chose que les tronqués \inserted{au cran 2}, dans la catégorie \uline{\unsure{homotopique}} des H-espaces - i.e. ce qu'on déduit des H-espaces par tuage des \unsure{groupes} d'homotopie en dimensions $\handwritten{>}2$.

\margin{\handwritten{\(\Vert\)}}% parallel strokes beside the question
(Demander à Quillen relations exactes entre \unsure{Objets|Hbjets|objets}-groupes et objets \deleted{ponctués} connexes de la catégorie homotopique \inserted{ponctuée}: forment-ils des catégories équivalentes par le foncteur espace des lacets\handwritten{?}).
% slanted
\margin{\handwritten{non, on a un foncteur \unsure{..|tr.|pl.} fidèle (?) mais \unsure{pas} \unsure{ess.t|ess.} surjectif...}}

Pour les \uline{vrais} groupes \inserted{abéliens} de $(2\text{-}\mathrm{Cat})$, l'approche de Quillen consiste à les identifier d'abord à des groupes abéliens de la catégorie $(\mathrm{Bicat})$, donc à des grpupes abéliens bi-semi\deleted{\unsure{e}}simpliciaux particuliers (conditions sur chaque ligne et chaque colonne). Appliquant la version bi-semisimpliciale du tapis de Dold, on trouve que \deleted{cette ca} la catégorie de ces objets est équivalente à la catégorie des \deleted{objets} bicomplexes \inserted{de chaines} de groupes \unsure{abéliens} qui sont nuls en dehors du carré $\deleted{\unsure{D}}0\handwritten{\leq}i,j\handwritten{\leq}1$. \deleted{Quillen dit que} La \deleted{colonne}\inserted{\handwritten{cond.}} que la première ligne corresponde à un \deleted{ensemble} groupe semisimplicial trivial signifie alors qu'elle n'a que le terme de degré zéro, i.e. que le bicomplexe est réduit à

% Arrow strokes in the bicomplex are in pen.
\diagram{\begin{tikzcd}[ampersand replacement=\&,row sep=small,column sep=small]
 C_{oo} \& 0\arrow[l] \& \handwritten{0}\arrow[l]\\
 \deleted{\unsure{C|0}}C_{1o}\arrow[u] \& C_{11}\arrow[l]\arrow[u] \& \handwritten{0}\arrow[l]\\
 \handwritten{\unsure{0}}\arrow[u] \& \&
\end{tikzcd}}

% archive sheet number: 8
\scanpage{10}\mspage{\handwritten{9}}

\deleted{\unsure{L}}Ces\deleted{\unsure{e}} bicomplexes s'identifient aux complexes de \unsure{càbines|cahines|chbines|chaines} de \unsure{loggueur|longueur} 2. Donc la catégorie de $(2\text{-}\mathrm{Cat})$-groupes abéliens est équivelente à celle des complexes de chaines de longueur 2. L'explicitation de la 2-catégorie associée à un tel complexe de chaines est mon interprétation chérie. Il faudrait encore expliciter \deleted{les morphismes} 2-Grab-catégorie des homomorphismes (en un sens à préciser) entre deux\deleted{\unsure{es}} telles 2-Grab-catégories strictes, en termes d'un complexe $\mathrm{RHom}$.

L'idée qui se dégage de cette approche, et que Quillen on espère va expliciter, c'est que la \deleted{notion de catégorie de} notions de groupe \uline{abélien} dans $(n\text{-}\mathrm{Cat})$ correspond à celle de complexe de \inserted{chaines de} longueur $n$, en raisonnant soit mod. isomorphisme de part et d'autre, soit mod. une notion naturelle de $n$-équivalence d'un côté, et dans la catégorie dérivée de l'autre. Ceci doit pouvoir se préciser en utilisant des $\mathrm{Hom}^{\bullet}$ resp. des $\mathrm{RHom}^{\bullet}$.

Si maintenant $C$ est une catégorie avec un ``produit tensoriel'' \unsure{associatif} (et éventuellement unitaire), i.e. un bifoncteur avec des isomorphismes d'associativité (mais sans inverse ...), on peut lui associer \deleted{\unsure{(}}sans doute une Gr-catégorie universelle (tout comme à un monoïde on associe \unsure{un} groupe symmétrisé), d'où des invariants, qu'on pourra désigner par $K_{\unsure{o|0}}(C)$ et $K_1(C)$. On suppose qu'on pourra de même trouver, pour tout $n$, une Gr-$n$-catégorie universelle où envoyer $C$ de façon compati\unsure{bme|ble} avec $\otimes$ (\unsure{tout|dout} cela dans un sens à préciser), dont les invariants $\handwritten{\pi}_n$ seraient par définition les $K_n(C)$. \unsure{En|Un} cas intéressant est celui où $C$ est une catégorie avec produits finis (par exemple une catégorie additive) et $\otimes$ est le produit; \deleted{mais} \deleted{(}\handwritten{[}Mais on notera que dans le cas d'une catégorie non seulement \unsure{additive} mais abélienne, il y a une notion plus adéquate du $K_o$, qui devrait correspondre à une notion plus adéquate de la suite des $K_i$, inspirée du tapis de Roos; sans doute c'est en termes de (``vraie'') catégories triangulées

% archive sheet number: 9
\scanpage{11}\mspage{\handwritten{10}}

qu'i\unsure{l|h} conviendra de poser les définitions adéquates). - Peut-être \unsure{Quillen} a-t-il une définition semi-si\unsure{mp|m}pliciale plus directe, sans passer par des $n$-catégories, des $K_i$ et plus précisément du H-espace qui devrait correspondre à $C$.\margin{\handwritten{\(\Vert\)\ \uuline{\unsure{; oui|si oui|oui}}}}
% slanted
\deleted{Dans le cas} Ce n'est certainement \inserted{en général} pas l'ensemble semi-simplicial $SC$ \deleted{engendré \unsure{à}} associé à $C$, puisque dans le cas d'une catégorie additive, $C$ est homotopiquement trivial (en effet le foncteur identique est ``\unsure{homotope}'' au foncteur nul, puisqu'il suffit d'avoir un homomorphisme d'un foncteur dans l'autre; d'où une équivalence d'homotopie entre $C$ et la catégorie nulle). Pourtant, il semble bien que dans le cas où $C$ est une \unsure{Gr-catégorie} i.e. admet des inverses, alors ce soit bien la construction ``naïve'' $SC$ qui est la bonne, puisqu'on trouve bien un $\handwritten{\pi}_{\unsure{o|0}}$ et un $\handwritten{\pi}_1$ qui sont aussi les $\handwritten{\pi}_o$ et $\handwritten{\pi}_1$ au sens catégorie sans Gr-structure.\margin{\handwritten{? \(\vert\)}}

\margin{\handwritten{\unsure{??|?} \(\vert\)}}\handwritten{À expliquer : pour des Gr-$n$-catégories, les $\pi_i$ au sens semi-simplicial sont aussi les $\pi_i$ au sens Gr-$n$-catégories, savoir les objets à isom près dans divers $\underline{\mathrm{Hom}}(1,1)$ ...... \deleted{Il y a aussi une interprétation de l'opération \unsure{du|des} $\unsure{\pi_1|\pi_0|\Omega\mathpunct{:}}$ \unsure{des|sur} \unsure{autres|lacets} en termes d'opération de shift géométrique évidente sur les $n$-catégories.....}}

% archive sheet number: 10
\scanpage{12}\mspage{\handwritten{11}}

III \quad Point de vue ``motivique'' en théorie du cobordisme.

Soit $C$ la catégorie des variétés différe\unsure{b|n|m}tiables \inserted{\handwritten{(compactes?)}} (pas \unsure{nécessairement} \unsure{orinhtables|orientables}), \inserted{les morphismes étant les classes d'homotopie d'applications continues.} Si $B$ est une catégorie, on s'intéresse aux couples $(F_{\bullet},F^{\bullet})$ d'un foncteur covariant et d'un foncteur contravariant de $C$ dans $B$, satisfaisant les conditions que pour tout $X\handwritten{\in}\mathrm{Ob}\ C$, on a $F_{\bullet}(X)=F^{\bullet}(X)$, et que si on a un produit fibré ordinaire

% Four diagonal arrows drawn in pen between typed objects; labels are typed.
\diagram{\begin{tikzcd}[ampersand replacement=\&,row sep=small,column sep=small]
 \& X'\arrow[dl,"g'"']\arrow[dr,"f'"] \& \\
 X\arrow[dr,"f"'] \& \&Y'\arrow[dl,"g"]\\
 \&Y \& {,}
\end{tikzcd}}

avec $f$ et $g$ transversaux, alors on a $g^{\bullet}f_{\bullet}=\unsure{f!|f'_{\bullet}}\,g'^{\bullet}$ \inserted{(on pose $f_{\bullet}=F_{\bullet}(f),f^{\bullet}=F^{\bullet}(\unsure{f})$\unsure{)}}. La question est de trouver une $B$ et un couple $(F_{\bullet},F^{\bullet})$ qui soient ``universels'' dans un sens évident. Quillen prouve que la construction suivante donne une telle solution.

Pour toute variété $X\handwritten{\in}\mathrm{ob}\ C$, on définit un anneau commutatif $B(X)$ (l'anneau de bordisme de $X$) en prenant comme éléments les \unsure{classes} \deleted{de} d'éléments de $C$ au dessus de $X$, pour la relation de cobordisme: $f':X'\handwritten{\to}X$ et $f'':X''\handwritten{\unsure{-|\to}}X$ sont équivalents, s'il existe une variété à bord $Z$, et un $f:Z\handwritten{\to}X$, tel que $\mathrm{bord}\ Z=X'\handwritten{\amalg}X''$ et $f'$ et $f''$ sont induits par $f$. La somme de $B(X)$ est la somme disjointe (c'est une loi de \unsure{groupe} \deleted{l'inverse étant} où tout élément est d'ordre \deleted{1}\handwritten{2} \deleted{ou 2}), le produit est donné par le produit fibré, en utilisant le théorème de transversalité de Thom pour bouger un peu $f''$, de sorte à le rendre transversal à $f'$.

\handwritten{\uline{NB} L'anneau est de plus \uline{gradué} par la dimension relative....}

Si on a $f:X\handwritten{\to}Y$, on définit $f_{\bullet}:B(X)\handwritten{\to}B(Y)$ de façon évidente (\unsure{toute} variété sur $X$ est sur $Y$), et $f^{\bullet}:B(Y)\handwritten{\to}B(X)$ par produit fibré (utilisant encore le th. de transversalité). De cette façon, $B(X)$ devient un foncteur-anneau \inserted{\handwritten{(\unsure{gradué})}} contravariant en $X$, et un foncteur groupe covariant en

% archive sheet number: 11
\scanpage{13}\mspage{\handwritten{12}}

$X$. De plus, on a la formule \inserted{\handwritten{du type $g^{\bullet}f_{\bullet}=f'_{\bullet}g'^{\bullet}$ ci-dessus, \unsure{impliquent|impliquant} la formule}} de projection habituelle. Ceci posé, les constructions bien connues en théorie des correspondances permettent de définir une catégorie $\underline{B}$ dont lds objets sont ceux de $C$, et les $\mathrm{Hom}$ étant donnés par
\[
\mathrm{Hom}_B(X,Y)=B(X\mathbin{x}Y)\ .
\]

On a un isomorphisme canonique $\underline{B}\handwritten{\simeq}\underline{B}^{o}$, qui est l'identité sur les objets et la symmétrie $B(X\mathbin{x}Y)\handwritten{\simeq}B(Y\mathbin{x}X)$ sur les flêches, et un foncteur évident \deleted{\unsure{$C\longrightarrow\underline B$}} $F_{\bullet}:C\handwritten{\to}\underline{B}$, qui est l'identité sur les objets et \unsure{donné} par les graphes sur les flêches. Composé avec la symmétrie précédente, on trouve un foncteur $F^{\bullet}:\deleted{\unsure{B(}}\,C\handwritten{\to}\underline{B}$. On vérifie enfin que le couple $(F_{\bullet},F^{\bullet})$ satisfait aux conditions envisagées dans l'alinéa 1 (ceci \unsure{aussi} devrait être formalisé au cas des correspondances générales ...). On vérifie enfin qu'il a la propriété universelle qu'il faut. Lorsqu'on se donne un couple $(F'_{\bullet},F'^{\bullet})$ à valeursdans une catégorie $B'$, il faut trouver une factorisation par $B$. Elle est toute trouvée sur les objets, il faut la définir sur les flêches, donc définir, pour $X,Y$ dans $C$, une application
\begingroup
\makeatletter\tagsleft@true\makeatother
\[
B(X\mathbin{x}Y)\handwritten{\longrightarrow}\mathrm{Hom}(X',Y')\ . \tag*{$(\ast)$}
\]
\endgroup
Mais si $Z$ est au\deleted{d}dessus de $X\mathbin{x}Y$, d'où $p:Z\handwritten{\to}X$ et $q:Z\handwritten{\to}Y$, on lui \unsure{associe} un élément du deuxième membre de $(\ast)$ comme le composé
\[
X'\xrightarrow{F'^{\bullet}(p)} Z'\xrightarrow{F'_{\bullet}(\deleted{\unsure{p}}q)}Y'\ .
\]

Il faut prouver que l'élément obtenu en remplaçant $Z$\inserted{\handwritten{\unsure{$Z_o$|$(=Z_0)$}}} par un $Z_1$ cobordant au dessus de $X$ \inserted{\handwritten{est le même}}. Or si $g:T\handwritten{\to}X\mathbin{x}Y$ réalise le cobordisme, on voit (regardant seulement $T$ et la décomposition $\mathrm{bord}(T)=Z_{\unsure{o|0}}\handwritten{\amalg}Z_1$) qu'il existe une variété $\handwritten{\overline{T}}$ sur $S^1$ (obtenue en ``doublant'' $T$) \deleted{et une application $T$ \unsure{sur}} trqnsversale au dessus d'un couple de points $s_o,\unsure{s_1}$ de $S^1$, tels que $Z_i$ soit l'image inverse \inserted{\handwritten{dans $\overline{T}$}} de $s_i$ ($i=0,1$); d'autre part,

% archive sheet number: 12
\scanpage{14}
\mspage{\handwritten{13}}

$g$ manifestement se prolonge en $\bar{g}:\bar{T}\longrightarrow X\mathbin{x}Y$. Cela nous ramène à \unsure{prouver} que \inserted{le composé} $\handwritten{\unsure{F'|F^i|F^o}}\bar T'\xrightarrow{\handwritten{\alpha_i}} \handwritten{\unsure{F'|F^i|F^o}}Z'_i\xrightarrow{\handwritten{\alpha_i}}\handwritten{F}\bar T'$ ne dépend pas de $i=0,1$. Ecrivons $X$ au lieu de $\bar T'$, \unsure{et} considérons le diagramme cartésien d'inclusions

% The square and its arrowheads are penned; the graph and alpha labels are handwritten.
\diagram{\begin{tikzcd}[ampersand replacement=\&]
X_s\arrow[r,"\handwritten{\alpha_s}"]\arrow[d,"\handwritten{\alpha_s}"'] \& X\arrow[d,"\handwritten{\Gamma\unsure{\triangleq|=|\underline{=}}(p,\mathrm{id}_X)}"]\\
X\arrow[r,"i_s"'] \& S^1\mathbin{x}X\mathpunct{,}
\end{tikzcd}}

on a alors $\handwritten{\alpha}{}_{s}^{\handwritten{*}}\handwritten{\alpha}{}_{s\handwritten{*}}=\Gamma^{\handwritten{*}}i_{s\handwritten{*}}$, qui ne dépend pas de $s$ puisque $i_s$ n'en dépend pas à homotopie près. - Le reste est sans doute l'AQT.

Sur la strcture des $B_{\bullet}(X)$. Ce sont tous des algèbres graduées sur $B_{\bullet}(\mathrm{pt})=\handwritten{\Lambda}$, qui est une algèbre graduée à degrés positifs, réduite à $\underline{F}_2$ en degré\deleted{\unsure{x|s}} $0$. C'est d'ailleurs une algèbre de polynômes connue\inserted{\handwritten{(à une infinité de générateurs, $X_i$ \unsure{$i$ entier|en tout|$i$ entier $\geq 1$} \unsure{$\ne$} \unsure{de la forme} $2$\illegible{1})}}.
\margin{\handwritten{\unsure{?|??}}}

On a un homomorphisme \inserted{d'annaàuxpgradués} évident $B_{\bullet}(X)\longrightarrow H^{\unsure{\mathbb B|B}}(X)$, nul en degrés $\handwritten{>}0$, qui à tout $Z$ sur $X$ associe l'image de $1\handwritten{\in} H^o(Z)$ par l'homomorphisme de \unsure{Gysin} associé à $Z\longrightarrow X$. Cet homomorphisme par Thom est surjectif (toute classe de cohomologie provient d'une image de variété compacte....); \deleted{\unsure{noter}}\unsure{.}
\margin{\handwritten{\uline{NB $H^*(X)$} $(=H^*(X,\mathbb F_2))$}}% slanted
\margin{\handwritten{en fait $H(X)\simeq B_{\bullet}(X)/\Lambda^+B_{\bullet}(X)$}}% slanted

On ppouve même qu'il existe un relèvement fonctoriel en $X$ $H^{\bullet}(X)\longrightarrow B_{\bullet}(X)$, \handwritten{(}\inserted{\handwritten{il est douteux qu'il soit}} compatible avec multiplication\handwritten{)} \deleted{\unsure{kin|kix}}tel que l'on ait:
\margin{\handwritten{?}}

\deleted{a)} L'homomorphisme correspondant $H^{\bullet}(X)\otimes_{\underline F_2}\handwritten{\Lambda\simeq} B_{\bullet}(X)$ est un isomorphisme (ce sera vrai pour tout relèvement individuellement ...), \deleted{et}

b\handwritten{)} \handwritten{[}\inserted{\handwritten{Quillen ignore si}} \deleted{\unsure{Le|Le ce}} relèvement est compatible aussi avec les hom. de Gysin\unsure{.}
\margin{\handwritten{?}}

Ainsi, la catégorie $\underline B$ admet un foncteur pleinement fidèle dans la catégorie des modules gradués libres de type fini sur $\handwritten{\Lambda}$, les modules de type fini qu'on obtient ainsi ayant des générateurs à degrés $\handwritten{<}0$ \unsure{et} satisfaisant (pour $X$ connexe i.e. quand il y a un seul générateur de degré zéro) la dualité de Poincaré. Quand on prend la catégorie \uline{\unsure{quasi} abélienne} $\inserted{\handwritten{\widetilde B}}$ associée à $\underline B$, on trouve donc tous les modules gradués libres de type fini à générateurs \deleted{négatifs} en degrés $\handwritten{\leq}0$ seulement. C'est un
\margin{\handwritten{i.e. l'enveloppe ``Karoubienne''}}% slanted

% archive sheet number: 13
\scanpage{15}
\mspage{\handwritten{14}}

catégorie \inserted{additive(et même quasi-abélienne)}\margin{\handwritten{$=$ karoubienne $+$ additive}} avec $\otimes$ qui a toutes les vertus, sauf d'être abélienne, et sauf d'admettre des \uline{Hom} i\deleted{\unsure{j|n}}nternes. Pour trouver ces derniers, on peut si on veut paraphraser la construction de la catégorie des motifs, \unsure{en} \deleted{introduisant un \unsure{motif}} \inserted{\handwritten{en}} considérant l'objet $\inserted{\handwritten{\Lambda}(-1)}$ de $\handwritten{\widetilde{\underline B}}$ conoyau de $p^{\handwritten{\unsure{*|\bullet}}}:e\longrightarrow S$, où $p:S^1\longrightarrow e$ est l'application du cercle sur un point. C'est un objet de poids $1$ (le poids est défini par le degré \deleted{\unsure{d}} des générateurs, changés de signe). On constate aussitôt que le foncteur $M\mapsto M(-1)=M\otimes\handwritten{\Lambda}(-1)$ de $\handwritten{\widetilde{\underline B}}$ dans lui-même est pleinement \unsure{fidèle}, ce qui permet alors, par une construction formelle standard, de comp\inserted{\handwritten{lé}}ter $\handwritten{\widetilde{\underline B}}$ en une catégorie où $\handwritten{\Lambda}(-1)$ devient inversible. Il se trouve alors que dans cette dernière les $\underline{\mathrm{Hom}}(X,Y)$ internes existent et \deleted{satisfont} sont isomorphes à $\overset{\handwritten{\vee}}X\otimes Y$, ... en fait , on trouve maintenant \handwritten{la} catégorie de\deleted{\unsure{x|s}} \uline{tous} les modules gradués libres de type fini sur $\handwritten{\Lambda}$. Bien sûr, cette catégorie n'est pas encore abélienne pour autant, puisque $\handwritten{\Lambda}$ n'est pas un corps; \deleted{pour}\inserted{si on veut} définir la catégorie abélienne engendrée, on s'attend à trouver la\deleted{s} catégorie desmodules de \deleted{type} \inserted{\handwritten{présentation}} \unsure{finie} sur $\handwritten{\Lambda}$. (\deleted{\unsure{---} est-il noethérien, i.e. le nombre de générateurs est-il fini ? ).}\inserted{\handwritten{qui est abélienne bien que $\Lambda$ ne soit pas \unsure{noethérien|noéthérien}, car $\Lambda$ est $\unsure{\underline{\lim}|\underset{\longrightarrow}{\lim}}$ d'anneaux \unsure{noéth.|noeth.} à hom de transition \uline{plats})}}.
\margin{\handwritten{\unsure{?|??}}}

Cette construction ressemblerait à celle d'une théorie des motifs où on n'aurait pas passé à l'équivalence homologique pour les cycles, mais où on aurait gardé l'équivalence rationelle, - qui est bien trop fine \deleted{\unsure{pou|pour}}à beaucoup d'égards. Pour trouver un équivalent plus étroit, il \unsure{convinnt|convint|convient|convinmt} de reprendre la théorie précédente,on pourrait formuler le même Pb universel que plus haut, mais en \unsure{prenant} l'équivalence homologique (via le graphe) des morphismes de variétés au lieu de l'équivalence homotopique. \deleted{Mais alors il semble que l'on trouve comme catégorie solution celle dont les objets sont les varié}

% archive sheet number: 14
\scanpage{16}
\mspage{\handwritten{15}}

\deleted{tés, les morphismes \uline{toutes} les correspondances \inserted{co}homologiques (en \unsure{vertu} du thorème de Thom rappelé plus haut). On trouve donc un foncteur pleinement fidèle de cette catégorie dans celle des vectoriels gradués \deleted{\unsure{de d|degr}} positivement de dim. finie sur $\underline F_2$, qui est déjà abélienne, et si on veut rendre comme plus haut $\underline F_2(-1)$ inversible, on \unsure{trouve} tous les vectoriels gradués de dim finie sur $\underline F_2$. C'est un peu \unsure{trop} trivial comme structure, et on peut faire mieux,} \deleted{en prenant comme correspondances \unsure{pour ceci|Pour ceci}}

\inserted{Pour trouver une catégorie universelle,} on associe à toute $X$ \deleted{le sous-anneau \unsure{ensemble} de} l'ensemble \inserted{$S(X)$} des sous-variétés \deleted{$S(X)$} de $X$, modulo équivalence homologique. Donc $S(X)\handwritten{\subset}H^{\bullet}(X)$, et on prouve qu'on trouve en fait un sous-anneau de $H^{\bullet}(X)$ (grâce à un théorème de Thom qui caractérise $S(X)$ \handwritten{en} termes d'opérations de Steenrod, cf plus bas). On prouve de même (par le même théorème) que les $S(X\mathbin{x}Y)$ sont stables par composition \unsure{de} correspondances cohomologiques\handwritten{(} cela ne semble pas trivial, puisque les $S(X)$ ne sont pas stables par Gysin...). D'où une catégorie additive dont les objets sont ceux de $C$, les flêches sont les \inserted{\handwritten{él. des}} $S(X\mathbin{x}Y)$, Il semble qu'elle doive jouer le rôle universel que j'ai envisagé plus haut. En tous cas, sa définition géo\deleted{\unsure{é}}métrique est fort simple Nous allons en donner une interpétation algébrique. Pŏur ce\unsure{ĉ|c}\unsure{î|i}, considérons l'anneau des opérations cohomologiques \deleted{additives sur \unsure{i.}} \inserted{``stables''} \deleted{des endomorphismes du \deleted{\unsure{le|la}} foncteur $X\longrightarrow H^{\bullet}(X)$ qui sont additifs, et qui en degré zéro se réduisent à un multiple de l'identité.} C'est un anneau gradué, avec application diagonale anticommutative\inserted{ i.e. commutative (car. $2$ !)}, associative et unitaire. Son dual gradué est une algèbre \inserted{graduée \unsure{antioommutative|anticommutative} donc \deleted{commutative},} commutative,avec application diagonale associative et unitaire; il est réduit à $\underline F_2$ en degré zéro. Donc il corespond à un schéma en groupes affine $G$ sur $\underline F_2$, qui est \unsure{d'ailleurs} le foncteur des $\otimes$-automorphismes \inserted{additifs} du foncteur $H^{\bullet}$ sur $C$ \inserted{\handwritten{qui \uline{opèrent} \unsure{trivialement} sur \deleted{\unsure{l'élément}} $\unsure{H^1(S^1)|H^1(\Delta^1)}$. \deleted{\unsure{On}}}}
\margin{\handwritten{\unsure{?|??} \(\vert\)}}% at the left of the definition of the stable operations

% archive sheet number: 15
\scanpage{17}
\mspage{\handwritten{16}}

Le théorème de Thom est alors qu'un $x\handwritten{\in}H^{\bullet}(X)$ est dans $S(X)$ sss il est \deleted{\unsure{stable|stabil}} \inserted{fixe} par $G$, i.e. satisfait $A^{\handwritten{+}}\mathbin{.}x=0$. (Pourquoi cette condition est-elle déjà nécessaire ?) U\inserted{\handwritten{n}}e autre façon de le voir est de dire que le foncteur $H^{\bullet}$ considéré comme allant de la catégorie $\underline S$ précédente dans la catégorie des \deleted{vectoriels} représentations de $G$ dans des vectoriels sur $\underline F_2$, est pleinement fidèle. \deleted{Comme $G$ est} Quand on prend l'enveloppe pseudo-abélienne, on trouve une sous-catégorie pleine de la catégorie des représentations de $G$, mais il est douteux qu'on trouve une catégorie abélienne, voire la catégorie de toutes les représentations de $G$ de \unsure{dimensio|dimension} finie sur $\underline F_2$. (Quelle est la structure de $G$, qui est prolisse; d'après Quillen, il serait pro-\unsure{unipotent|umipotent} ....). Ici $G$ bien sûr joue le rôle d'un groupe de Galois motivique, mais contrairement à la situation familière, il semble qu'il soit fortement non réductif (et d'ailleurs le corps de \unsure{défimtion|défintion|définition} dans tout ceci est $\underline F_2$, et non $\underline Q$ ! ).
\margin{\handwritten{\unsure{Oui}, car son \unsure{anneau} affine \unsure{a|est} une \unsure{infinité|suite|système|algèbre} de \unsure{générateurs|polynômes} ! une \unsure{quantité|infinité} d'\deleted{\unsure{indéte}}\unsure{indéterminées}, \deleted{\unsure{limite}} \unsure{sous} \unsure{relatio.|inductive|réunion} d'\unsure{algèbres} \unsure{polynômiales} à \unsure{augmentation|engendrement} \unsure{et} stables par $\Delta$}}% slanted

Quillen\deleted{t} donne une généralisation de la construction d'une catégorie universelle $\underline B$ au cas où on travaille avec des variétés pas nécessairement propres (mais $f_{\bullet}$ n'étant supposé défini que pour $f$ propre), et quand on introduit une condition d'orientation \inserted{stable} sur les \uline{morphismes} envisagés , via un morphisme $H\unsure{\text{---}|\longrightarrow}\mathrm{BO}(\infty)$, où $H$ est un H-espace (commutatif, associatif et unitaire ?): une orientation de $f:X\unsure{\text{---}|\longrightarrow}Y$ sera un relèvement homotopique de $X\longrightarrow\mathrm{BO}$, défini par \deleted{$\unsure{f^{\bullet}(T_X)|f^{*}(T_X)}$} $T_X-f^{\bullet}(T_Y)$, en $X\longrightarrow H$. les éléments de $H(X,Y)$ seront définis par des diagrammes de morphismes orientés $Z\longrightarrow X,\ Z\unsure{\text{---}|\longrightarrow}Y$, le deuxième étant propre, modulo bords de tels. En tant qu'ensemble, il se définit aussi en termes de la vari\handwritten{é}té $X\mathbin{x}Y$ munie simplement de la famille des supports formée des ensembles fermés propres sur $Y$ .....

\handwritten{Les $B^H_{\unsure{m|n}}(X)$ ($X$ compact disons) peuvent se définir homotopiquement à la Thom, comme les groupes d'homotopie stables d'\unsure{espaces|appl.} $\underline{\mathrm{Hom}}(X,\mathrm M\mathrm H(N-\unsure{m|n}))$ -----.}

\margin{\handwritten{cf exposé Karoubi à \uline{Bourbaki}}}

% archive sheet number: 16
